\section{Introducción}


En el estudio de máquinas eléctricas, los núcleos ferromagnéticos desempeñan un papel fundamental debido a su alta permeabilidad magnética, la cual permite confinar y conducir eficientemente el flujo magnético. Sin embargo, cuando estos núcleos son sometidos a campos magnéticos variables en el tiempo, se presentan fenómenos complejos como la histéresis y pérdidas de energía asociadas tanto al lazo de histéresis como a las corrientes parásitas (o de Foucault). 

La presente práctica de laboratorio tiene como objetivo analizar y comparar el comportamiento magnético de dos tipos de núcleos: uno macizo y otro laminado, mediante la visualización directa del ciclo de histéresis en un osciloscopio. Para lograr esto, se implementa un circuito integrador $RC$ que permite obtener una tensión proporcional a la densidad de flujo ($B$), mientras que una resistencia shunt se encarga de censar la corriente del devanado primario, la cual es directamente proporcional a la intensidad de campo magnético ($H$).

A partir de las mediciones experimentales realizadas, se calcularán las escalas correspondientes del ciclo, las pérdidas totales del núcleo, las pérdidas por histéresis (determinadas a partir del área encerrada por el ciclo) y las pérdidas por corrientes parásitas. La comparación cuantitativa entre ambos núcleos permitirá evidenciar la reducción de las corrientes parásitas en el núcleo laminado, validando así la importancia del uso de chapas magnéticas aisladas para mejorar la eficiencia energética en aplicaciones de corriente alterna. Finalmente, todos los resultados obtenidos se presentarán con sus correspondientes errores de medición y empleando las cifras significativas adecuadas.